% !TeX root = ../papers/nc-p3.tex
\section{有限点グラフの閉極限となる量子窓は二類だけである}
\label{sec:finite-quantum-closed-limits}
\begin{proposition}\label{prop:maximal-tangent-closed-limits}
三次元接点を持つ標準正則有限窓は半安定である。三生成基底 $B$ の点三次曲線が滑らかなら、
唯一の多重安定極限は $B[Z]$ であり、量子窓ではない。特異なら、位数2と4の二型についてそれぞれ、
中心文字 $Z$ と $xy=\zeta yx,yv=\zeta vy,vx=\zeta xv$ を持つ量子窓に退化する。
ここで $\zeta=-1$ 又は $\zeta^2=-1$ で、逆数の選択は同じ添字保持類を与える。
\end{proposition}
\begin{proof}
定理\ref{thm:maximal-tangent-ore}の移送表示で補文字を $Z$ とする。$Z$ より三つの基底文字を
高く取れば混合補正が消えて $B[Z]$ になる。この極限の関係は六本で、
四テンソルは三生成基底の巡回三体積 $w_B$ に $Z$ を交代符号で四通りに挿入したものである。
実際六関係の先頭部分はこれで、各辺の像は直和因子として極限と可換である。
三次体積の巡回性は滑らかな場合はHesse表示、特異な場合は定理\ref{thm:ternary-torsion-smoothing}の三関係表示による。
四テンソルの極限もこれらの辺の交叉の一意な直線に入る。

滑らかな場合の三体積は三文字の巡回置換と $(1,\omega,\omega^2)$ 倍で不変である（$\omega$ は原始三乗根）。
複素数上でこの二つのユニタリ行列は三文字上既約であり、各位置のGram行列の三文字部分はスカラーである。
$Z$ と三文字の成分間は、残る三位置の $Z$ の個数が異なり直交する。体積は四位置で巡回し、
全語に $Z$ がちょうど一個なので各位置でその重みのノルム平方は全体の四分の一である。
従ってGram行列は四文字全体でスカラーとなり、Kempf--Ness判定で閉軌道である。
標数零体上へは係数の有限生成定義体から移る。
この窓の点対応は $Z=0$ で滑らかな平面三次曲線を含み、その外で三文字が非零なら
混合の可換関係が平行移動の固定点を強制する。従って一次元支持は滑らかな三次曲線である。
量子窓の点支持は座標線型空間の和であるから同型にならない。

特異な三型では三生成基底の関係は巡回量子交換子に $0,1,2$ 個の純平方を加えたものである。
平方のある文字を残る文字より高く取ると、その各平方だけが初期関係で消える。
従って $B[Z]$ は記載した量子点へ落ちる。量子点は多重安定
（定理\ref{thm:quantum-face-resonance}の証明）であり、二閉類は面値で区別される。
\end{proof}

\begin{theorem}\label{thm:finite-quantum-closed-limits}
二隣接切断が階数6で四辺の点対応が有限な窓の軌道閉包に量子点があれば、
その量子点の四面値はすべて同じ $\zeta$ で、$\zeta=-1$ 又は $\zeta^2=-1$ である。
これら二つの添字保持類はいずれも実現する。特に多項式環のGITファイバーには
有限点グラフの窓は一つもない。元の窓に群対称性もAS性も仮定しない。
\end{theorem}
\begin{proof}[三つの必要な一次係数]
定理\ref{thm:quantum-universal-chambers}により元の窓は正則で、その三余指標の一つの正の変位で表示できる。
各辺の階数6は還元で保たれる。量子二項式 $r_{ij}$ ($i<j$) の各々を、それより真に高い重みだけで補正し、
六つの順序対 $ij$ での係数が単位行列である関係基底を取る。
背景で同じ射影文字を各位置に使える。余指標 $(0,1,4,6)$ では最初の二文字の両直線の積に一式しか残らず有限でない。
$(0,6,8,9)$ では上三文字どうしの三式は最低文字を含まないので最低点に三次元接点が生じる。
これは命題\ref{prop:finite-polynomial-limit-flag}のJacobi階数の議論と同じで、初期二項式の交換倍率は非零で足りる。
この場合の量子極限は直前の命題で尽きる。

$(0,3,5,9)$ では順に
\[
  r_{02}\longrightarrow 11,\qquad r_{12}\longrightarrow30,\qquad
  r_{03}\longrightarrow22
\]
への補正はすべて重み差1である。この順に係数を $A,B,C$ とする。
$A=0$ なら最初の二文字に $r_{01}$ 以外の制限がない。
$C=0$ なら最初の三文字に $r_{01},r_{02},r_{12}$ 以外の制限がない。
どちらも有限性に反する。$B=0$ なら $r_{12}$ に最低文字がなく、
$r_{13},r_{23}$ にもないので上と同じ三次元接点を得る。
一辺でもこの場合に入れば直前の命題を使える。
残りでは三係数が非零である。余重み1は初期の一次変形に直接現れる。
各補正の負添字対は $02,12,03$ であり、定理\ref{thm:quantum-face-resonance}の証明の
外積複体による必要条件は単位面や三重、四重の面値でも適用できる。
これらは文字変更の一次境界を持たない。従って
$h_0=h_2=h_1$ かつ $h_0=h_3$ である。面値の積は1なので共通値は四乗根である。
\end{proof}

\begin{proof}[共通値1での二次障害]
全て1の背景は通常の交代体積 $w$ である。上記余重みでのテンソルの最初の二項を
$w+V_1+V_2$ とし、$V_j$ の語の重みは $17+j$ とする。
二切断で $w$ の行列を $M$、$V_j$ の行列を $M_j$ と書く。
$P=(01,02,03,12,13,23)$、$I=M_{PP}^{-1}$、$L,R$ をそれぞれ左核の行、右核の列の基底とすれば、
Schur補の最初の二式は
\[
   L M_1 R=0,\qquad L M_2 R=L M_{1,*,P}\, I\,M_{1,P,*} R .
\]
両切断の式を積み重ねる。重み1では14列の線型部分の階数は11で、三自由変数を $a,b,c$ と取ると
上記の $(A,B,C)=(c,-b,-a)$ である。重み2では12列の線型部分の階数は7で、
右辺をその像で割った条件は $4ab=0$ となる。
\path{computations/verify_finite_quantum_closed_limits.py} が体積の24係数から両行列を構成し、
主元による消去と三つの関係係数を有理数上で検算する。
従って三次元接点がない場合の $BC\ne0$ と両立しない。
三次元接点があっても上述の閉極限は多項式環でない。
逆に二つの非自明型の背景は定理\ref{thm:quantum-tangent-finite}の
有限条件を満たす補正を選べば実現する。全ての面値が等しい場合に奇置換は
共通値を反転するので、残る類数は二つである。
\end{proof}
